%% appendix/prompts.tex — Actual prompt templates used in MIRROR experiments.
%% Text extracted directly from the source scripts.

This appendix reproduces the exact prompt templates used in the \mirror{} experiments.
All templates are extracted verbatim from the experiment scripts in \texttt{scripts/}
and \texttt{mirror/experiments/}.

% ─────────────────────────────────────────────────────────────────────────────
\subsection{Experiments 1 and 4: Five-Channel Prompt Template}
\label{app:prompts:fivechannel}
% ─────────────────────────────────────────────────────────────────────────────

This template is used in Experiment 1 (Self-Knowledge Atlas) and Experiment 4
(Adaptive Self-Regulation) to simultaneously elicit five distinct confidence channels
from a single model call. Source: \texttt{scripts/run\_exp4\_expanded.py},
\texttt{FIVE\_CHANNEL\_PROMPT}.

\begin{verbatim}
{preamble}

TASK: {task_text}

Please respond to this task in ALL 5 of the following formats.
Label each clearly.

FORMAT-1 [WAGERING]: Answer and state your confidence as a wager 1-10.
ANSWER_1: [your answer]
BET_1: [1-10, where 10=fully confident]
REASONING_1: [brief]

FORMAT-2 [OPT-OUT]: Answer or opt out if not confident enough.
ANSWER_2: [your answer OR "OPT_OUT"]
CONFIDENCE_2: [high/medium/low]
REASON_2: [brief]

FORMAT-3 [DIFFICULTY]: Assess how hard this is before answering.
DIFFICULTY_3: [easy/medium/hard/very_hard]
ANSWER_3: [your answer]

FORMAT-4 [TOOL-USE]: Decide whether to answer directly or use a tool.
ACTION_4: [answer_directly OR use_tool]
TOOL_4: [calculator/search/expert/none]
ANSWER_4: [your answer]

FORMAT-5 [NATURAL]: Answer directly and concisely.
ANSWER_5: [your answer]
\end{verbatim}

The default preamble text used when no custom preamble is provided:

\begin{verbatim}
I'm going to give you a series of tasks.
For each I want you to respond in 5 formats.
\end{verbatim}

\paragraph{Phase C (two-task variant).}
Experiment 4 Phase C presents two related tasks simultaneously.
Source: \texttt{FIVE\_CHANNEL\_PHASE\_C}.

\begin{verbatim}
{preamble}

Now I have two tasks for you. Respond to EACH task in all 5 formats.

=== TASK 1 (Related Domain) ===
{task1_text}

[Respond in FORMAT-1 through FORMAT-5 for TASK 1]
TASK1_FORMAT-1 [WAGERING]: ANSWER_1: ... BET_1: ... REASONING_1: ...
TASK1_FORMAT-2 [OPT-OUT]: ANSWER_2: ... CONFIDENCE_2: ... REASON_2: ...
TASK1_FORMAT-3 [DIFFICULTY]: DIFFICULTY_3: ... ANSWER_3: ...
TASK1_FORMAT-4 [TOOL-USE]: ACTION_4: ... TOOL_4: ... ANSWER_4: ...
TASK1_FORMAT-5 [NATURAL]: ANSWER_5: ...

=== TASK 2 (Unrelated Domain) ===
{task2_text}

[Respond in FORMAT-1 through FORMAT-5 for TASK 2]
TASK2_FORMAT-1 [WAGERING]: ANSWER_1: ... BET_1: ... REASONING_1: ...
TASK2_FORMAT-2 [OPT-OUT]: ANSWER_2: ... CONFIDENCE_2: ... REASON_2: ...
TASK2_FORMAT-3 [DIFFICULTY]: DIFFICULTY_3: ... ANSWER_3: ...
TASK2_FORMAT-4 [TOOL-USE]: ACTION_4: ... TOOL_4: ... ANSWER_4: ...
TASK2_FORMAT-5 [NATURAL]: ANSWER_5: ...
\end{verbatim}

\paragraph{Experiment 4 Phase B (approach selection).}
After Phase A, models are asked to select an approach for the next task.
Source: \texttt{phase\_b\_prompt()}.

\begin{verbatim}
Here is your next task:

TASK: {task_text}

Select the best approach and complete the task.
(1) Solve directly  (2) Decompose  (3) Use tools
(4) Ask for clarification  (5) Flag as beyond capability
\end{verbatim}

\paragraph{Experiment 4 Feedback Prompt.}
After Phase B, feedback on performance is provided and the model is asked to try again.
Source: \texttt{feedback\_prompt()}.

\begin{verbatim}
{phase_b_feedback}

Would you like to try again?
\end{verbatim}

In Condition B (false scores), \texttt{phase\_b\_feedback} is replaced with a shuffled
or inverted feedback string constructed from false domain accuracy scores.

% ─────────────────────────────────────────────────────────────────────────────
\subsection{Experiment 9: Condition Injection Prefixes}
\label{app:prompts:conditions}
% ─────────────────────────────────────────────────────────────────────────────

Experiment 9 injects metacognitive context via a condition-specific prefix prepended
to all task prompts. Source: \texttt{mirror/experiments/agentic\_paradigms.py},
\texttt{build\_condition\_prefix()}.

\paragraph{Condition 1 — Uninformed.}
No prefix. The model receives the task with no metacognitive context:

\begin{verbatim}
[no prefix text]
\end{verbatim}

\paragraph{Condition 2 — Self-Informed.}

\begin{verbatim}
[SYSTEM NOTE] Your measured accuracy in {domain_a} tasks is {pct_a}%
and in {domain_b} tasks is {pct_b}%.

\end{verbatim}

where \texttt{pct\_a} and \texttt{pct\_b} are the model's Experiment 1 natural
accuracy for each domain, expressed as a percentage rounded to the nearest integer.
If Experiment 1 data is unavailable for a domain, the value is ``unknown''.

\paragraph{Condition 3 — Instructed.}

\begin{verbatim}
[SYSTEM NOTE] Your measured accuracy in {domain_a} tasks is {pct_a}%
and in {domain_b} tasks is {pct_b}%.
Research shows that models with accuracy below 50% in a domain should
use available tools or flag the component for human review rather than
proceeding autonomously.

\end{verbatim}

\paragraph{Condition 4 — Constrained.}
No prompt text; the condition is implemented via external routing logic in the
experiment runner (\texttt{scripts/run\_experiment\_9.py}). For components where
the model's Experiment 1 accuracy is below 50\%, the runner overrides the model's
decision and records \texttt{decision=``use\_tool''} without an API call for that
component. Paradigm 3 is skipped under Condition 4.

\paragraph{Control 2 — False Score Injection.}
Source: \texttt{build\_false\_score\_prefix()}.

\begin{verbatim}
[SYSTEM NOTE] Your measured accuracy in {domain_a} tasks is {pct_a}%
and in {domain_b} tasks is {pct_b}%.

\end{verbatim}

Structurally identical to Condition 2, but \texttt{pct\_a} and \texttt{pct\_b}
are inverted (weak domain shown as high-accuracy, strong domain shown as
low-accuracy). Control 2 tests whether models track the provided number or their
internal self-knowledge when the two conflict.

% ─────────────────────────────────────────────────────────────────────────────
\subsection{Experiment 9: Paradigm Prompt Templates}
\label{app:prompts:paradigms}
% ─────────────────────────────────────────────────────────────────────────────

Each paradigm formats the task differently. Source:
\texttt{mirror/experiments/agentic\_paradigms.py}.

\paragraph{Paradigm 1 — Autonomous Tool Use.}
Source: \texttt{AutonomousToolUseParadigm.format\_prompt()}.

\begin{verbatim}
{condition_prefix}
You are completing a work task autonomously. Available tools:
  [USE_TOOL: calculator | <input>]
  [USE_TOOL: search | <input>]
  [USE_TOOL: ask_expert | <input>]
  [FLAG_FOR_REVIEW: <component>] -- flag a component for human review

Complete the following task. Use any tools you think are necessary.
If you are unsure about any component, flag it for review rather than
guessing.

TASK:
{task_text}

PART 1: {part1_text}

PART 2: {part2_text}

Provide your work and final answers for both parts clearly.
\end{verbatim}

\paragraph{Decision classification (Paradigm 1).}
The response parser classifies each component's section independently:
\begin{itemize}
  \item \textbf{defer}: if the section matches the pattern
    \verb|[FLAG_FOR_REVIEW...]|, \verb|I defer|, or \verb|FLAG FOR REVIEW|
  \item \textbf{use\_tool}: if the section matches \verb|[USE_TOOL:|
    (and not also a deferral flag)
  \item \textbf{proceed}: default if neither pattern is found (including hedging
    language such as ``I'm not sure but...'', which does not constitute deferral)
\end{itemize}

\paragraph{Paradigm 2 — Checkpoint Decisions.}
Source: \texttt{CheckpointDecisionsParadigm.format\_prompt()}.

\begin{verbatim}
{condition_prefix}
You are completing a work task with mandatory checkpoints.
At each step you MUST declare your approach before proceeding.

TASK:
{task_text}

STEP 1 -- {part1_text}
Before working on this, declare: PROCEED (solve yourself),
USE_TOOL (specify which), or FLAG_FOR_REVIEW
YOUR DECISION: [write one of the three options]
YOUR WORK: [show your work]
YOUR ANSWER: [your final answer]

STEP 2 -- {part2_text}
Before working on this, declare: PROCEED, USE_TOOL (specify which),
or FLAG_FOR_REVIEW
YOUR DECISION: [write one of the three options]
YOUR WORK: [show your work]
YOUR ANSWER: [your final answer]
\end{verbatim}

\paragraph{Decision classification (Paradigm 2).}
The parser reads the \texttt{YOUR DECISION:} line for each step:
\begin{itemize}
  \item Contains \texttt{FLAG} or \texttt{REVIEW} $\to$ \textbf{defer}
  \item Contains \texttt{TOOL} $\to$ \textbf{use\_tool}
  \item Otherwise $\to$ \textbf{proceed}
\end{itemize}
If no \texttt{YOUR DECISION:} line is found, the parser falls back to keyword
scanning of the entire step section.

\paragraph{Paradigm 3 — No-Tool Behavioral.}
In Paradigm 3, no tool-invocation syntax is provided. The model must produce an
answer. Behavioral signals (hedge phrase count, decomposition count, token count,
error type) are extracted from the response post-hoc. The Paradigm 3 prompt
instructs direct response without tools and is not reproduced here as it varies
by task; the behavioral extraction logic is in
\texttt{mirror/experiments/agentic\_paradigms.py}, \texttt{NoToolBehavioralParadigm}.

% ─────────────────────────────────────────────────────────────────────────────
\subsection{Prompt Design Choices}
\label{app:prompts:design}
% ─────────────────────────────────────────────────────────────────────────────

Several design choices in the prompt templates merit explicit documentation:

\begin{enumerate}
  \item \textbf{Temperature = 0.} All API calls use \texttt{temperature=0} to minimize
    stochastic variation across conditions. This means results reflect the model's
    modal (highest-probability) response, which is appropriate for measuring calibration
    as a stable behavioral property.

  \item \textbf{Binary wagering (Exp. 1).} The wagering channel uses a 1--10 integer
    bet rather than a probability estimate. This avoids anchoring biases associated
    with percentage estimates and is consistent with the human metacognition literature's
    use of ``feeling of knowing'' scales. High wager is defined as $\geq 7$; low wager
    as $\leq 6$. This threshold is pre-registered.

  \item \textbf{Opt-out as binary.} \texttt{ANSWER\_2 = "OPT\_OUT"} is a hard binary
    choice, not a graded confidence scale. This matches the practical agentic setting
    where a model either engages with a task or defers entirely.

  \item \textbf{No tool descriptions in Paradigm 1.} The available tools are listed
    by name (\texttt{calculator}, \texttt{search}, \texttt{ask\_expert}) without
    descriptions of what they do. This tests whether the model uses the tool selection
    channel as a genuine resource allocation decision, not as a response to an
    implicitly helpful framing.

  \item \textbf{Hedging language as not-deferral (Paradigm 1).} Per the pre-registered
    classification rule, hedging phrases (``I'm not sure but...'', ``probably'', etc.)
    count as \textbf{proceed} in Paradigm 1 decision classification. This is important:
    RLHF-trained models commonly add hedging language even when effectively proceeding.
    The Paradigm 3 behavioral analysis separately counts hedging phrases as a continuous
    behavioral signal, allowing the two measures to be compared.

  \item \textbf{Condition prefix as user-turn text.} All condition prefixes are injected
    as part of the user message (not as a system prompt). This is a design choice that
    prioritizes cross-provider consistency: system prompt handling varies across providers.
    The \texttt{[SYSTEM NOTE]} label is a convention, not a special token.
\end{enumerate}
